\begin{table}[H]
\centering
\caption{Welfare Losses from Demand vs.\ Supply Shocks}
\begin{threeparttable}
\small
\begin{tabular}{lcccc}
\toprule
& \multicolumn{4}{c}{CE Welfare Loss (\%)} \\
\cmidrule(lr){2-5}
Scenario & Risk-neutral & $\sigma=1$ & $\sigma=2$ & $\sigma=5$ \\
\midrule
\multicolumn{5}{l}{\textit{Panel A: Demand shock (Great Recession analog)}} \\[3pt]
Baseline & 0.92 & 0.92 & 0.92 & 0.89 \\
No scarring ($\lambda = 0$) & 0.84 & 0.84 & 0.83 & 0.81 \\
No OLF exit ($\chi = 0$) & 0.93 & 0.93 & 0.93 & 0.91 \\
\midrule
\multicolumn{5}{l}{\textit{Panel B: Supply shock (COVID analog)}} \\[3pt]
Baseline & 0.08 & 0.0808 & 0.0827 & 0.0888 \\
\midrule
Demand/Supply ratio & 12 & 11 & 11 & 10 \\
\bottomrule
\end{tabular}
\begin{tablenotes}[flushleft]
\small
\item \textit{Notes:} Consumption-equivalent (CE) welfare losses under different risk preferences. Risk-neutral: linear utility $u(c) = c$. CRRA: $u(c) = c^{1-\sigma}/(1-\sigma)$; $\sigma = 1$ is log utility, $\sigma = 2$ is standard macro calibration, $\sigma = 5$ is high risk aversion. The demand shock reduces aggregate productivity by 5\% following an AR(1) process $a_t = a_{ss}(1 - 0.05 \cdot \rho_a^t)$. The supply shock doubles the separation rate for 3 months. Welfare is computed over a 600-month simulation horizon with terminal-value continuation.
\end{tablenotes}
\end{threeparttable}
\label{tab:welfare}
\end{table}
